% ==================== Research Notes Preamble ====================
% Document Type: Research Notes
% Compiler: LuaLaTeX
% Fonts: Latin Modern (LaTeX default)
% =================================================================

% --- Layout ---
\usepackage[left=1.25in, right=1.25in, top=1in, bottom=1in]{geometry}

% --- Math ---
\usepackage{amsmath}
\usepackage{amsthm}

% --- Color ---
\usepackage[dvipsnames]{xcolor}
\definecolor{hex191919}{RGB}{25, 25, 25} % body text (off-black)
% Link colour: one blue for citations, cross-references, and URLs (matches the paper).
\definecolor{hexblue}{RGB}{0, 0, 180}
\color{hex191919}

% --- Graphics / Tables ---
\usepackage{graphicx}
\usepackage[export]{adjustbox}
\usepackage{booktabs}
\usepackage{threeparttable}
\usepackage{tabularx}
\usepackage{array}
\usepackage{rotating}

% --- Math / Symbols ---
\usepackage{mathtools}
\usepackage{dsfont}
\usepackage{latexsym}
\usepackage{cancel}
\usepackage[thinc]{esdiff}

\newtheorem{proposition}{Proposition}
\newtheorem{definition}{Definition}

% --- Citations / Bibliography ---
\usepackage[round,comma]{natbib}
\setlength{\bibsep}{0em}
\AtBeginDocument{\setlength{\bibhang}{\parindent}} % reference hang = paragraph indent
\bibliographystyle{aer}
\bibpunct{(}{)}{,}{a}{}{;}

% --- Hyperlinks ---
\usepackage{xr-hyper}
\usepackage[
  backref=page,
  colorlinks=true, % coloured link text, no boxes
  linkcolor=hexblue,
  citecolor=hexblue,
  urlcolor=hexblue,
  filecolor=hexblue
]{hyperref}
\renewcommand*{\backref}[1]{}
\renewcommand*{\backrefalt}[4]{%
  \ifcase #1 (Not cited.)%
  \or [#2]%
  \else [#2]%
  \fi
}

% --- Typography enhancements ---
\usepackage[expansion=false]{microtype}  % Improves text appearance with subtle spacing adjustments

% --- Captions ---
\usepackage{caption}
\usepackage{subcaption}
\captionsetup{labelsep=period, font=normalsize, labelfont=normalsize}
\captionsetup[subfigure]{labelformat=simple, labelsep=period, font=normalsize, labelfont=normalsize}
\captionsetup[subtable]{labelformat=simple, labelsep=period, font=normalsize, labelfont=normalsize}
% Spacing inside floats: one gap, \housefloatgap (0.5em), between every caption or
% subcaption and its content, between panel rows, and before the Notes. It matches the
% gap between reference-list entries, so a float reads as one tight unit.
\newlength{\housefloatgap}\setlength{\housefloatgap}{0.5em}
\captionsetup{skip=\housefloatgap}
\captionsetup[sub]{skip=\housefloatgap}

% --- TOC formatting ---
\usepackage[titles]{tocloft}
\renewcommand{\cftdot}{}
\cftpagenumbersoff{subsection}
\cftpagenumbersoff{subsubsection}
\setlength{\cftbeforetoctitleskip}{0pt}
\setlength{\cftaftertoctitleskip}{0pt}

% --- Section heading formatting ---
\usepackage[noindentafter]{titlesec}
\titlelabel{\thetitle.\quad}

\titleformat{\section}[hang]{\normalsize\bfseries\centering}{\textup{\thesection.}}{0.5em}{}[]
\titlespacing{\section}{0em}{1em}{1em}

\titleformat{\subsection}[hang]{\normalsize\itshape\centering}{\textup{\thesubsection.}}{0.5em}{}[]
\titlespacing{\subsection}{0em}{1em}{1em}

\titleformat{\subsubsection}[hang]{\normalsize\itshape}{\textup{\thesubsubsection.}}{0.5em}{}[]
\titlespacing{\subsubsection}{0em}{1em}{1em}

\titleformat*{\paragraph}{\normalsize\bfseries}
\titlespacing{\paragraph}{\parindent}{1em}{0.5em}

\titleformat*{\subparagraph}{\normalsize\itshape}
\titlespacing{\subparagraph}{\parindent}{1em}{0.5em}

% --- Spacing / Paragraphs ---
\usepackage{setspace}
\renewcommand{\baselinestretch}{1}
\usepackage{indentfirst}
% One indent for the whole document, 1em (the base unit of the spacing system):
% paragraphs, run-in \paragraph headings, footnote first lines (via \bodyparindent),
% and the hanging indent of reference-list entries (\bibhang, set to match below).
\setlength{\parindent}{1em}

% --- Footnotes ---
\usepackage[bottom]{footmisc}
\renewcommand{\hangfootparindent}{0pt}
\renewcommand{\hangfootparskip}{0pt}
\renewcommand{\footnotemargin}{0pt}
\renewcommand{\footnotelayout}{\hspace{0pt}}
\newlength{\bodyparindent}
\AtBeginDocument{\setlength{\bodyparindent}{\parindent}}

% No footnote rule: the space above (\skip\footins) separates footnotes from the text
\renewcommand{\footnoterule}{}
\setlength{\skip\footins}{0.25in}

% Footnote text formatting
\usepackage{etoolbox}
\makeatletter
\patchcmd\@makefntext{\hss\@makefnmark}{\@makefnmark\hss}{}{\fail}
% First line indented by the body's \parindent (saved below at full size, so the mark
% lines up with paragraph starts in the text even though footnotes are set smaller);
% wrapped lines return to the margin. Matches the body, whose paragraphs all indent.
\long\def\@makefntext#1{\leavevmode
  \hskip\bodyparindent
  \@makefnmark\nobreak
  \hskip1pt\relax#1%
}
\makeatother

% --- Listings / Subfiles / Standalone ---
\usepackage{listings}
\usepackage{subfiles}
\usepackage{standalone}

% --- Date formatting ---
\usepackage{datetime}
\newdateformat{monthyeardate}{\monthname[\THEMONTH] \THEYEAR}

% --- Appendix ---
\usepackage[toc,page]{appendix}
\renewcommand{\appendixpagename}{\centering\Large{\MakeUppercase{\textbf{Appendix}}}}
\renewcommand{\appendixtocname}{Appendix}

% --- Subcaption numbering ---
\renewcommand{\thesubfigure}{\Alph{subfigure}}
\renewcommand{\thesubtable}{\Alph{subtable}}

% --- Custom title / abstract ---
\makeatletter
% Title block: a lighter version of the paper's (normal-size title, no top space),
% with the same em-based gaps between lines.
\renewcommand{\maketitle}{\bgroup\setlength{\parindent}{0pt}
  \centering
  \normalsize\MakeUppercase{\@title}\\
  \vspace{1em}
  {\@author}\\
  \vspace{1em}
  {\@date}\\
  \vspace{1em}
\egroup
}
\makeatother

\renewenvironment{abstract}
 {\normalsize
 \itshape
  \list{}{
    \setlength{\leftmargin}{0.375in}%
    \setlength{\rightmargin}{\leftmargin}%
  }%
  \item\relax}
 {\endlist}

% --- Helpers ---
\newcommand{\specialcell}[2][c]{%
  \begin{tabular}[#1]{@{}c@{}}#2\end{tabular}}

\newcommand\blfootnote[1]{%
  \begingroup
  \renewcommand\thefootnote{}\footnote{#1}%
  \addtocounter{footnote}{-1}%
  \endgroup
}

% --- Float tweaks ---
\makeatletter
\setlength{\@fptop}{0pt}
\makeatother

% --- Header/footer (loaded but not activated) ---
\usepackage{fancyhdr}

% --- TOC hide/unhide helpers ---
\newcounter{oldtocdepth}
\newcommand{\hidefromtoc}{%
  \setcounter{oldtocdepth}{\value{tocdepth}}%
  \addtocontents{toc}{\protect\setcounter{tocdepth}{-10}}%
}
\newcommand{\unhidefromtoc}{%
  \addtocontents{toc}{\protect\setcounter{tocdepth}{\value{oldtocdepth}}}%
}

\makeatletter
\renewcommand\tableofcontents{%
  \@starttoc{toc}%
}
\makeatother

% --- House table & float style (copy this block to paper/slides/notes) ---
% Rules: uniform width, air around each rule, flush double rule at top.
% \sym is esttab's star macro (undefined in fragment mode). Needs booktabs + array.
\providecommand{\sym}[1]{\ensuremath{^{#1}}}
\setlength{\heavyrulewidth}{0.4pt}
\setlength{\lightrulewidth}{0.4pt}
\setlength{\aboverulesep}{0.28em} % = the former 0.65ex, now in em
\setlength{\belowrulesep}{0.28em}
% flush double rule: two lines with only \doublerulesep between (no inflated gap)
\newcommand{\doublerule}{%
  \specialrule{\arrayrulewidth}{\aboverulesep}{0pt}%
  \specialrule{\arrayrulewidth}{\doublerulesep}{\belowrulesep}}
% Use right after a heading when a float immediately follows it (no text between):
% cancels only the float's top \intextsep, so it doesn't hang below the title.
\newcommand{\tightafterheading}{\par\vspace{-\intextsep}}
% Float spacing: 2em between a float and body text (top/bottom of page or in-text) and
% between stacked floats -- the "separate block" step of the spacing system (0.5em inside
% a float, 1em around headings, 2em around floats). In em, so the paper (1.3-spaced) and
% notes (single-spaced) get the same gaps. A little stretch helps page breaks.
\setlength{\floatsep}{2em plus 0.5em minus 0.5em}
\setlength{\textfloatsep}{2em plus 0.5em minus 0.5em}
\setlength{\intextsep}{2em plus 0.5em minus 0.5em}
\setlength{\belowcaptionskip}{0pt}

% --- House float commands (copy this block to notes; slides place floats manually) ---
% Panels are given as a body of \panel commands (NOT a delimited string), so calls read
% top-to-bottom. Each panel: \panel[sub=<subcaption>, label=<sublabel>]{path}. Both keys
% optional -- omit sub= for a plain (unlettered) panel; omit label= for no \label.
% Trailing args are always {caption}{label}{notes}; label/notes omitted if empty; notes
% span \textwidth. Paths root-anchored (../../output/...). Starred form adds
% \tightafterheading (first float under a heading, no text between). Non-regular floats
% (\ContinuedFloat, sideways) stay hand-written -- see fig-background.tex.
% Depends on: etoolbox (\ifstrempty), keyval (loaded by caption), subcaption.
%
%   \housefigure[<w>|cont]{ \panel{...} \panel{...} ... }{caption}{label}{notes}
%       <w> = uniform width of EVERY panel (and the lone figure if there's one panel).
%       Default: \textwidth for a single panel, 0.5\textwidth for 2+. Rows wrap by how
%       many <w>-wide panels fit across \textwidth (so [\textwidth]=1/row, [.5\tw]=2/row,
%       [.33\tw]=3/row); a short last row centers. A \panel with sub= gets an (A)/(B)
%       subcaption; with none, it is a bare image -- so an all-plain call renders as a
%       simple single/grid figure. \housesubfigure is a kept alias of \housefigure.
%       [cont] (the literal word) makes this a \ContinuedFloat continuation of the
%       PRECEDING \housefigure (same figure, no new number): caption becomes a bare
%       \caption* (label arg ignored) and the panel body may be empty for a
%       footnote-only continuation page -- pair with an intervening \clearpage, e.g.
%       \housefigure{...}{Title}{lbl}{}  \clearpage  \housefigure[cont]{}{Title (cont'd)}{}{notes}
%   \housetable[cont]{ \panel{frag.tex} ... }{caption}{label}{notes}
%       One \panel = one \input booktabs fragment wrapped in a subtable; multiple panels
%       stack vertically, each with an "A." subcaption header when sub= is given (bare
%       \input if not). label= makes the panel cross-referenceable like a subfigure:
%       \subref{lbl} -> "A", \ref/\autoref -> the panel number+letter. [cont] (the
%       literal word) makes this a \ContinuedFloat continuation of the PRECEDING
%       \housetable (same table, no new number): caption becomes a bare \caption*
%       (label arg ignored) and the panel body may be empty for a footnote-only
%       continuation page -- pair with an intervening \clearpage, e.g.
%       \housetable{...}{Title}{lbl}{}  \clearpage  \housetable[cont]{}{Title (cont'd)}{}{notes}
\makeatletter
\newcounter{housefign}\newcounter{housefigi}\newcounter{housefigcol}%
\newcounter{housefigperrow}\newcounter{housetabi}%
\newlength{\housefig@w}\newlength{\housefig@acc}%
\define@key{housepanel}{sub}{\def\housepanel@sub{#1}}%
\define@key{housepanel}{label}{\def\housepanel@label{#1}}%
% \panel is only meaningful inside a house float; error if used loose.
\newcommand{\panel}[2][]{\PackageError{housefloat}{\string\panel\space is only valid
  inside \string\housefigure\space or \string\housetable}{}}
% Pass 1 (in a throwaway box): count panels, to pick the default width.
\newcommand{\housefig@countpanel}[2][]{\stepcounter{housefign}}
% Compute how many \housefig@w-wide panels fit across \textwidth (>=1).
\newcommand{\housefig@computeperrow}{%
  \setcounter{housefigperrow}{1}\setlength{\housefig@acc}{2\housefig@w}%
  \loop\ifdim\housefig@acc<\dimexpr\textwidth+0.01\textwidth\relax
    \stepcounter{housefigperrow}\addtolength{\housefig@acc}{\housefig@w}%
  \repeat}
% Pass 2: render one figure panel.
\newcommand{\housefig@figpanel}[2][]{%
  \def\housepanel@sub{}\def\housepanel@label{}\setkeys{housepanel}{#1}%
  \stepcounter{housefigi}%
  % \relax (not a space) terminates each number scan: a trailing space would be KEPT
  % when the branch is true, injecting ~3.33pt between panels and wrapping a full row.
  \ifnum\value{housefigi}=1\relax
    \setcounter{housefigcol}{1}%
  \else
    \ifnum\value{housefigcol}<\value{housefigperrow}\relax
      \hfill\stepcounter{housefigcol}%
    \else
      \par\nointerlineskip\vskip\housefloatgap\setcounter{housefigcol}{1}%
    \fi
  \fi
  \begin{subfigure}[t]{\housefig@w}\centering
    % \ifdefempty (tests the macro's EXPANSION), not \ifstrempty (which would see the
    % single macro token as non-empty and always fire, printing a bare subcaption letter).
    % \expandafter expands \housepanel@sub ONCE so \caption records the actual text: else
    % subcaption writes the literal token \housepanel@sub to the .aux, and a doc that
    % \externaldocument's this one (undefined there) dies with "Undefined control sequence".
    \ifdefempty{\housepanel@sub}{}{\expandafter\caption\expandafter{\housepanel@sub}}%
    \ifdefempty{\housepanel@label}{}{\label{\housepanel@label}}%
    \includegraphics[width=\linewidth]{#2}%
  \end{subfigure}%
  \ignorespaces}
\newcommand{\housefigure}{\@ifstar{\tightafterheading\housefigure@}{\housefigure@}}
\newcommand{\housefigure@}[5][]{%
  \ifstrequal{#1}{cont}{\def\housefig@cont{1}}{\def\housefig@cont{}}%
  \setcounter{housefign}{0}\sbox\z@{\let\panel\housefig@countpanel#2}%
  \ifdefempty{\housefig@cont}{%
    \ifstrempty{#1}{%
      \ifnum\value{housefign}=1
        \setlength{\housefig@w}{\textwidth}\else\setlength{\housefig@w}{0.5\textwidth}\fi
    }{\setlength{\housefig@w}{#1}}%
  }{\setlength{\housefig@w}{\textwidth}}%
  \housefig@computeperrow
  \begin{figure}[ht!]
    \centering
    \ifdefempty{\housefig@cont}{\caption{#3}\ifstrempty{#4}{}{\label{#4}}}{\ContinuedFloat\caption*{#3}}%
    \setcounter{housefigi}{0}\setcounter{housefigcol}{0}%
    \let\panel\housefig@figpanel\ignorespaces#2%
    \ifstrempty{#5}{}{\par\nointerlineskip\vskip\housefloatgap
      \begin{minipage}{\textwidth}{\footnotesize\textit{Notes.} #5\par}\end{minipage}}%
  \end{figure}%
}
\let\housesubfigure\housefigure  % consolidated: subfigures are just panels with sub=
% Render one table panel as a REAL subtable (parallels \housefig@figpanel): panels
% stack vertically, each with an "A." caption when sub= is given. A real subtable (not
% a plain-text header) is what makes \ref/\subref/\autoref resolve to the panel letter
% (A/B) rather than the outer table number. A caption-less subtable (no sub=) prints
% nothing and centers identically, so a single bare-panel table is unchanged.
\newcommand{\housetab@panel}[2][]{%
  \def\housepanel@sub{}\def\housepanel@label{}\setkeys{housepanel}{#1}%
  \stepcounter{housetabi}%
  \ifnum\value{housetabi}=1\relax\else\par\nointerlineskip\vskip\housefloatgap\fi
  \begin{subtable}{\textwidth}\centering
    % \expandafter so \caption records the text, not the macro token (see
    % \housefig@figpanel); the "A." + period come from \captionsetup[subtable].
    \ifdefempty{\housepanel@sub}{}{\expandafter\caption\expandafter{\housepanel@sub}}%
    \ifdefempty{\housepanel@label}{}{\label{\housepanel@label}}%
    \input{#2}%
  \end{subtable}\par}
\newcommand{\housetable}{\@ifstar{\tightafterheading\housetable@}{\housetable@}}
\newcommand{\housetable@}[5][]{%
  \begin{table}[ht!]
    \centering
    \ifstrempty{#1}{\caption{#3}\ifstrempty{#4}{}{\label{#4}}}{\ContinuedFloat\caption*{#3}}%
    \setcounter{housetabi}{0}\let\panel\housetab@panel#2%
    \ifstrempty{#5}{}{\par\nointerlineskip\vskip\housefloatgap
      \begin{minipage}{\textwidth}{\footnotesize\textit{Notes.} #5\par}\end{minipage}}%
  \end{table}%
}
\makeatother

% ==================== END PREAMBLE ====================
